\chapter{Introducción}
\label{chap:introduccion}

Crear sistemas, automatizar tareas y expresar ideas de un dominio mediante
programación puede suponer una barrera importante para los usuarios principiantes.
La literatura sobre programación basada en bloques muestra que estos entornos
se utilizan a menudo para introducir a usuarios noveles en la
programación, porque ofrecen una entrada más accesible que los lenguajes
textuales tradicionales \cite{lin2021landscape}. Esta dificultad no aparece
exclusivamente en la enseñanza de la informática. También afecta a usuarios
no expertos que podrían utilizar la programación para crear simulaciones,
automatizar procesos o representar conocimiento de un dominio.

Los lenguajes basados en bloques reducen parte de esta dificultad. En lugar de
escribir instrucciones textuales desde cero, el usuario trabaja con piezas
visuales que se arrastran y se conectan. La forma de las piezas, sus colores, sus
etiquetas y la organización por categorías ayudan a reconocer qué construcciones
son posibles \cite{weintrop2018modalities}. Herramientas como Scratch, App
Inventor, Tynker o MakeCode muestran esta idea en contextos educativos y en
dominios como robótica, ciencia de datos, desarrollo móvil o creación de juegos
\cite{lin2021landscape}.

Blockly es una de las bibliotecas más usadas para crear editores de este tipo.
Su documentación lo describe como una biblioteca web que permite integrar en una
aplicación un editor visual basado en bloques. Blockly se ocupa de la
representación visual, el arrastre, el encaje de piezas y la gestión del
workspace. El desarrollador, en cambio, debe definir los bloques, los campos, las
conexiones y la generación asociada \cite{blocklyWhatIs,blocklyDocs}.

Por lo tanto, Blockly facilita la creación de editores visuales, pero no suprime
todo el esfuerzo necesario. Cada nuevo dominio sigue necesitando tipos de bloque, campos,
entradas, categorías de toolbox, reglas de conexión, serialización,
validaciones y, en algunos casos, generación de código. Este esfuerzo debe
repetirse cuando se quiere crear un editor de bloques para otro dominio. Además, los
entornos de programación por bloques no siguen un único diseño. La literatura
indica que cada entorno toma decisiones propias sobre la organización de los
bloques, la relación con el código textual y la transición entre programación
visual y textual \cite{lin2021landscape,weintrop2018modalities}.

Esta situación plantea un problema de ingeniería de lenguajes. Si los conceptos
principales de un dominio ya están descritos en un modelo o metamodelo, no
debería ser necesario duplicar la misma información como código
JavaScript de Blockly. La ingeniería dirigida por modelos propone precisamente
trabajar con modelos como artefactos principales para especificar, diseñar,
probar y generar software \cite{brambilla2017modelDriven}. También se han señalado
motivaciones como la reutilización, la reducción del tiempo de desarrollo y la
mejora de calidad \cite{hutchinson2014industry,akdur2018survey}. Desde esta perspectiva, un editor
Blockly puede verse como una sintaxis visual generada a partir de una
descripción abstracta del dominio.

El ecosistema Eclipse ofrece tecnologías adecuadas para explorar esta idea. EMF
permite definir modelos estructurados y generar código a partir de ellos. Su
metamodelo Ecore representa clases, atributos, referencias, herencia y
cardinalidades \cite{emfOverview,steinberg2008emf}. Xtext permite definir
lenguajes textuales de dominio específico y genera infraestructura de parsing,
enlace, validación y edición \cite{xtextProject,bettini2013xtext}. Estas
tecnologías permiten describir dominios con precisión, pero no generan
automáticamente una sintaxis visual basada en bloques.

El objetivo de este Trabajo Fin de Máster es desarrollar un entorno para crear
lenguajes de dominio específico con sintaxis basada en bloques. La idea no es
construir un único editor de bloques, sino una infraestructura que permita
describir un dominio y generar de forma automática el código necesario para
editarlo con Blockly.

La herramienta desarrollada se llama Model2Blockly. Se entrega como un conjunto
de plugins Eclipse y como una cadena de generación HTML/JavaScript. Admite dos
rutas de entrada. La primera parte de metamodelos Ecore, que pueden enriquecerse
con anotaciones sobre la sintaxis visual. En esta ruta, las clases se convierten
en tipos de bloque, los atributos en campos, las referencias de contención en
entradas estructurales y las cardinalidades en reglas de comprobación. La
segunda ruta parte de una sintaxis textual \texttt{.m2b} definida con Xtext. Esa
sintaxis es un DSL de definición de lenguajes: permite al diseñador del
lenguaje declarar qué bloques, campos, conexiones, categorías y validaciones
tendrá el editor Blockly generado.

Ambas rutas se transforman en un modelo intermedio común,
\texttt{Blockly\allowbreak{}Editor\allowbreak{}Spec}. Este modelo contiene los
elementos necesarios para generar el editor: tipos de bloque, categorías, campos,
entradas de valor, referencias, validaciones y metadatos de interfaz. La
instancia intermedia se serializa como XMI y se recarga antes de producir los
artefactos finales. A partir de esta representación se genera un editor Blockly
en HTML/JavaScript que puede ejecutarse directamente en el navegador, junto con
un informe de generación, un modelo de ejemplo y una vista visual de las reglas
de validación.

El desarrollo se ha organizado de forma incremental. Primero se generaron bloques
y paletas a partir de modelos de dominio. Después se añadieron validaciones para
detectar campos obligatorios vacíos o bloques requeridos que faltan. También se
incorporaron referencias dinámicas mediante listas desplegables, opciones de
personalización visual, anotaciones de interfaz, referencias opuestas,
validaciones semánticas simples, generación de código y exportación del modelo
creado por el usuario como JSON o XMI de instancia. La herramienta se evalúa
mediante AppMaker, un ejemplo conductor descrito por las dos rutas: como
metamodelo Ecore anotado y como fichero textual \texttt{.m2b}. Esta doble
descripción permite comparar la autoría en Ecore y en el DSL textual, y mostrar
que ambas entradas se normalizan hacia el mismo modelo intermedio antes de
generar el editor Blockly.

Las aportaciones principales del TFM son las siguientes. En primer lugar, se
define una arquitectura que conecta Ecore, un modelo intermedio EMF y Blockly
mediante transformaciones explícitas y serialización XMI. En segundo lugar, se
implementa el adaptador Ecore que extrae clases, atributos, referencias,
cardinalidades y anotaciones. En tercer lugar, se implementa una ruta textual
\texttt{.m2b} con Xtext y un adaptador que transforma ese modelo al mismo
\texttt{EditorSpec}. En cuarto lugar, se desarrolla un generador de editores
Blockly con bloques, toolbox, páginas HTML, serialización, validaciones,
exportación JSON/XMI e informes de generación. En quinto lugar, se añaden
referencias, personalización visual, generación de código y una representación
editable de reglas de validación que puede aplicarse de vuelta a la fuente. Por
último, se evalúa la solución con un ejemplo conductor, pruebas automáticas y
artefactos desplegados públicamente.

El alcance del trabajo se centra en la generación de editores Blockly desde
descripciones de dominio de alto nivel, ya sean metamodelos Ecore anotados o
modelos textuales \texttt{.m2b}, y en la reutilización de la arquitectura
mediante un modelo intermedio común. Quedan fuera del objetivo principal la importación
genérica de modelos XMI existentes, la integración completa de restricciones OCL
arbitrarias y la creación de interfaces finales muy especializadas para cada
dominio. La contribución central es reducir el esfuerzo necesario para crear
editores Blockly específicos mediante metamodelos, DSL textual, transformaciones
de modelos y generación de código.

El resto del documento se organiza de la siguiente manera. El capítulo 2
presenta los objetivos, el alcance y los criterios de evaluación del trabajo. El
capítulo 3 revisa el estado del arte y las tecnologías en las que se apoya la
propuesta. El capítulo 4 describe el diseño y la arquitectura de la solución,
mientras que el capítulo 5 detalla la implementación de Model2Blockly. El
capítulo 6 expone el ejemplo conductor y la evaluación de la herramienta. Por
último, el capítulo 7 recoge las conclusiones, las limitaciones identificadas y
las posibles líneas de trabajo futuro.
